\chapter{Theory}
\section{Superconductivity Fundamentals}



\subsection{Basic Phenomenology}
Most electrical conductors follow a clear temperature trend: as temperature decreases, \ac{DC} resistance decreases. This behavior stems from the mechanisms of conduction in metals, where two primary scattering processes limit conductivity. The first is electron-phonon scattering: as electrons move through the metal lattice, they collide with thermally vibrating ions, and this scattering reduces conductivity. The second is scattering from lattice imperfections. Because a real metal never has a completely crystalline lattice—always containing defects such as impurities, vacancies, or dislocations—these imperfections provide an additional source of scattering that limits conductivity. \cite{AshcroftSolidState}

Examining the temperature dependence of each mechanism shows that phonon scattering depends on temperature, since lattice vibrations intensify as temperature rises. Defect scattering, by contrast, is largely temperature-independent because static structural imperfections do not change substantially with temperature. Accordingly, one might expect a metal's conductivity to increase as temperature falls, leveling off once phonon scattering becomes negligible and conduction is limited solely by lattice imperfections. In reality, however, many metals deviate from this expectation: their electrical resistivity drops abruptly to zero below a characteristic transition temperature, known as the \ac{Tc}.
\subsection{Type-II Superconductors}
\subsection{Critical Surface}

\subsection{\textit{Thin Film Superconductors?}}



\section{Quenches in Superconductors}
\subsection{Quench Origins}
\subsection{Quench Dynamics}
\subsection{Consequences of undetected Quenches}
\subsection{Quench Protection}

\section{Quench Detection Methods}
\subsection{Voltage based Quench Detection}
\subsection{Alternative Quench Detection Methods}
\subsection{Shortcomings of Existing Methods}

\section{Superconducting Quench Detection}
\ac{SQD}
\subsection{Principle of Operation}
\subsection{Materials}
\subsection{Literature Examples}
\subsection{Design}

\section{Experimental Setup}